% SPDX-License-Identifier: LPPL-1.3c
% Copyright (C) 2026 Christophe Jorssen
% Author and Current Maintainer: Christophe Jorssen <christophe.jorssen@gmail.com>
% LPPL maintenance status: maintained
% This file is part of luafemm. See LICENSE and LICENSES.md for its terms.
\documentclass[varwidth=150mm,border=3mm]{standalone}
\usepackage{tikz,pgfplots}
\usetikzlibrary{femm}
\usepgfplotslibrary{femm}
\pgfplotsset{compat=1.18}
\newbox\ResponseCalculation
\tikzset{femm/materials/machine iron/.style={mu r=300},
  machine winding/.style={femm/region={material=copper,circuit=drive,turns=#1}},
  response field/.style={femm/lines=19,femm/arrows=false,
    femm/field style={blue!65!black}}}
\pgfkeys{/femm/measure/rotor load/.style={quantity=force,method=band}}
\begin{document}
\centering
% Frequency belongs to the template. Start coarse while arranging the figure.
\foreach \Frequency in {0,50}{%
  \setbox\ResponseCalculation=\hbox{%
    \begin{tikzpicture}[femm/problem={xmin=-10,xmax=10,ymin=-10,ymax=10,
        mesh size=1,depth=10,frequency=\Frequency},
      femm/circuit={name=drive,current=10},
      femm/cache={file=\jobname-template-\Frequency,mode=auto},
      femm/sliding band={name=machine-\Frequency,
        inner radius=3.5,outer radius=4.5,segments=64,layers=2}]
      \path[femm/region={material=machine iron,mesh size=.5}]
        (0,0) ellipse[x radius=3,y radius=1.5];
      \path[machine winding=100] (-7,-2) rectangle (-6,2);
      \path[machine winding=-100] (6,-2) rectangle (7,2);
    \end{tikzpicture}}}
% Capture the converged DC state at this mechanical angle.
\begin{tikzpicture}[femm/problem={sliding band=machine-0,rotor angle=30},
    femm/cache={file=\jobname-bias,mode=auto},
    femm/operating point={name=machine-bias},scale=2.5]
  \pic {femm geometry};
  \pic[response field] {femm field};
  \femmmeasure[rotor load,name=dc-load]
  \draw[teal,densely dashed] (0,0) circle[radius=4];
  \node[font=\small,fill=white] at (0,9) {DC bias: $I_0=10\,\mathrm A$};
\end{tikzpicture}\hspace{5mm}%
% The response inherits this exact mesh and angle; it changes only excitation.
\begin{tikzpicture}[femm/problem={operating point=machine-bias,frequency=50,
    circuit currents={drive=1},circuit phases={drive=30}},
    femm/cache={file=\jobname-response,mode=auto},scale=2.5]
  \pic {femm geometry};
  \pic[response field,femm/phase=0] {femm field};
  \femmmeasure[rotor load,name=small-load]
  \node[font=\small,fill=white] at (0,9) {$\widehat{\delta I}=1e^{i30^\circ}\,\mathrm A$};
\end{tikzpicture}
% An independent AC calculation has no DC bias and gives mean + 2f torque.
\setbox\ResponseCalculation=\hbox{%
  \begin{tikzpicture}[femm/problem={sliding band=machine-50,rotor angle=30,
      circuit phases={drive=30}},femm/cache={file=\jobname-ac,mode=auto}]
    \femmmeasure[rotor load,name=ac-load]
  \end{tikzpicture}}
\par\medskip
\begin{tikzpicture}
\begin{axis}[width=6.3cm,height=4.8cm,font=\small,xmin=0,xmax=360,xtick={0,90,180,270,360},
    grid=major,xlabel={Excitation phase (degrees)},ylabel={$T$ (mN\,m)},
    y filter/.expression={1000*y},scaled y ticks=false,title={AC without bias}]
  \femmforceplot[result=ac-load,component=torque,
    plot options={blue!65!black,thick,mark=none}]
\end{axis}
\end{tikzpicture}%
\begin{tikzpicture}
\begin{axis}[width=6.3cm,height=4.8cm,font=\small,xmin=0,xmax=360,xtick={0,90,180,270,360},
    grid=major,xlabel={Excitation phase (degrees)},ylabel={$\delta T$ (mN\,m)},
    y filter/.expression={1000*y},scaled y ticks=false,title={Response around DC bias}]
  \femmforceplot[result=small-load,component=torque,
    plot options={orange!80!black,thick,mark=none}]
\end{axis}
\end{tikzpicture}
\par\smallskip
\small Fixed rotor angle $30^\circ$, linear illustrative iron, depth $10\,\mathrm{mm}$.
The right-hand waveform is the first-order change, excluding the DC torque.
\end{document}
